\section{Literature Study}

\subsection{Human Activity Recognition}
Activity recognizers take RGB, skeleton, IMU, radar, or Wi-Fi, and usually
return one label from a closed set. An open-set model can reject an unknown
action. It does not describe that action. An open-vocabulary model can match
a clip to a class name while the action is still going on. Free-form
captions also exist. OV-HAR turns a sensor window into a sentence
\cite{ray2025ovhar}, and OpenHOUSE calls a vision-language model for a
free-form description of streaming video \cite{kang2025openhouse}. Both are
evaluated on recorded data, not on a live phone link. The missing piece is a
sentence produced on the phone, before any server round trip, while the
person is still moving.
\subsection{Vision-Language Models}
CLIP retrieves a sentence from a fixed bank by cosine similarity. It does
not write a new sentence \cite{radford2021learning}. A generative
vision-language model writes the sentence token by token, and accuracy
follows model size. Qwen2-VL scales this family up to 72B parameters
\cite{wang2024qwen2vl}. Models that can sit on a phone also exist.
MiniCPM-V is in the 2B--8B range \cite{yao2024minicpm}, and SmolVLM-256M
fits in under 1~GB \cite{marafioti2025smolvlm}. Published scores for these
models cover documents, charts, and general questions, not fine differences
between similar activities. A small model can write the sentence. It does
not, by itself, say which of those sentences are safe to leave on the
screen.
\subsection{Client--Server Collaboration}
Neurosurgeon profiles network layers and chooses where to split them, so
the server receives image features and the user waits for that transfer
\cite{kang2017neurosurgeon}. Speculative decoding splits the work by
tokens: a small model drafts, a large model verifies, and the server still
reads the picture \cite{leviathan2023fast}. edgeVLM answers on the device
at once and reuses a late cloud caption as context for later frames, rather
than as the result of the frame that is on screen now
\cite{qian2025edgevlm}. EdgeBoost uploads an input only when the local
confidence is low. On image classification this removes about 20--55\% of
the communication of a cloud-only run, and the cloud still returns a class
label \cite{said2024edgeboost}. What these systems do not do is show an
activity sentence immediately and send the picture only when that sentence
is unsafe.
\subsection{Communication}
A proportional-fair scheduler and an L4S queue assign service to packets
that have already been formed. They do not look at how long the laptop will
spend on the picture after it arrives
\cite{jalali2000data,briscoe2023l4s}. Semantic communication scores a
transmission by task success rather than by throughput, but the payload is
still a feature tensor \cite{xie2021deep}. Offloading chooses local or
remote execution, and adaptive streaming changes how many bytes are sent.
After the sentence has already been written, these mechanisms still do not
decide whether the picture should follow, how large that picture should be,
and which phone should be checked first.
Our system takes that last decision. Token confidence decides whether the
JPEG is uploaded. JPEG quality and resolution decide its size. For every
upload, the predicted edge latency sets the uplink deadline, and the
application-layer scheduler serves the earliest deadline first.
